\SGASection{7}{Relative root datum}
\SGARefTarget{sga3:xxvi:sec:7}

Throughout this section, \(S\) denotes a nonempty connected
semilocal scheme, \(G\) an \(S\)-reductive group, \(Q\) a maximal
split subtorus of \(G\), and \(L\) the centralizer of \(Q\) in \(G\):
\[
 L=\Centr_G(Q).
\]

\SGASubsection{7.1}{}
\SGARefTarget{sga3:xxvi:subsec:7:1}

Since \(Q\) is the largest central split subtorus of \(L\), every
element of \(G(S)\) which normalizes \(L\) also normalizes \(Q\).
Thus \textup{(cf.\ \SGARefLink{sga3:xxvi:subsec:7:1:1}{7.1.1})}
\[
 \Norm_G(L)(S)=\Norm_G(Q)(S).
\]
On the other hand, \SGARefLink{sga3:xxvi:corollary:6-4}{6.4} shows that
\[
 G(S)\longrightarrow(G/L)(S)
\]
is surjective.  Consequently there is a canonical identification
\[
 W_G(Q)(S)
 =\bigl(\Norm_G(Q)/\Centr_G(Q)\bigr)(S)
 \simeq\Norm_G(L)(S)/L(S).
\]

Let
\[
 M=\operatorname{Hom}_{S\text{-}\mathrm{gr}}
      (Q,\mathbf G_{m,S}),
\]
so that \(Q\simeq D_S(M)\) canonically.  Denote by \(W\) the group
of automorphisms of \(M\) defined by \(W_G(Q)(S)\).  Thus
\[
 W\simeq W_G(Q)(S)\simeq\Norm_G(L)(S)/L(S).
\]

\SGASubsection{7.1.1}{}
\SGARefTarget{sga3:xxvi:subsec:7:1:1}

In general one does not have
\(\Norm_G(L)=\Norm_G(Q)\).  For example, let \(S\) be the spectrum of
a field \(k\) having a quadratic extension \(k'\), and take for
\(G\) the unitary group
\[
 {}^2\!A_{2,k'/k}
\]
\textup{(cf.\ \SGARefLink{sga3:xxiv:3.11.2}{Exposé XXIV, 3.11.2})}.  The minimal parabolic
subgroups of \(G\) are its Borel subgroups, and their Levi subgroups
are maximal tori; hence
\[
 \bigl(\Norm_G(L)/L\bigr)(k)=\mathfrak S_3.
\]
Since \(G\) is not split, its maximal split tori have dimension at
most \(1\), and are therefore isomorphic to \(\mathbf G_{m,k}\).
The group \(\Norm_G(Q)/L\) acts faithfully on \(Q\), and
\[
 \bigl(\Norm_G(Q)/L\bigr)(k)=\mathbf Z/2\mathbf Z.
\]
\footnote{\emph{Editor's note.} The French edition prints
\(\Norm_G(Q)/Q\) in this final equality, although the immediately
preceding action and the definition of the relative Weyl group use
\(\Norm_G(Q)/L\).}

\SGASubsection{7.2}{}
\SGARefTarget{sga3:xxvi:subsec:7:2}

Let \(P\) be a parabolic subgroup of \(G\) with Levi subgroup \(L\);
such a subgroup exists by \SGARefLink{sga3:xxvi:cor:6:2}{6.2}.  It is necessarily minimal, by \SGARefLink{sga3:xxvi:rem:6:18}{6.18}.
Using the conjugacy of the minimal parabolic subgroups of \(G\)
\textup{(\SGARefLink{sga3:xxvi:corollary:5-7}{5.7})}, the conjugacy of the Levi subgroups of a parabolic
subgroup \textup{(\SGARefLink{sga3:xxvi:cor:1:8}{1.8})}, and the equalities
\[
 P=\Norm_G(P),\qquad \Norm_G(L)\cap P=L
\]
from \SGARefLink{sga3:xxvi:proposition:1-6}{1.6}, we obtain the following conclusion: the set of
\textup{(minimal)} parabolic subgroups of \(G\) having Levi subgroup
\(L\) is a principal homogeneous space under \(W\).

\SGASubsection{7.3}{}
\SGARefTarget{sga3:xxvi:subsec:7:3}

The Lie algebra of \(G\) decomposes under the action of \(Q\) as
\[
 \operatorname{Lie}(G)
 =\operatorname{Lie}(L)
  \oplus\bigoplus_{\alpha\in R}\operatorname{Lie}(G)^\alpha,
\]
where \(R\) is the set of nonzero characters \(\alpha\) of \(Q\)
such that \(\operatorname{Lie}(G)^\alpha\ne0\).  These are the roots
of \(G\) relative to \(Q\).

Let
\[
 M^*=\operatorname{Hom}_{S\text{-}\mathrm{gr}}
       (\mathbf G_{m,S},Q).
\]
It is in duality with \(M\), and \(W\) acts on it naturally by
transport of structure.

\begin{sgastatement}{\SGAInlineRefTarget{sga3:xxvi:theorem:7-4}Theorem 7.4}
With the notation of \SGARefLink{sga3:xxvi:subsec:7:3}{7.3}, there is a unique map
\[
 R\longrightarrow M^*,\qquad \alpha\longmapsto\alpha^*,
\]
which makes \((M,M^*)\) into a root datum \textup{(\SGARefLink{sga3:xxi:subsection:1-1}{Exposé XXI,
1.1})} with Weyl group \(W\).

Moreover, parabolic subgroups \(P\) of \(G\) having Levi subgroup
\(L\) correspond bijectively to systems of positive roots \(R^+\)
in \(R\) by the relation
\[
 \operatorname{Lie}(P)
 =\operatorname{Lie}(L)
  \oplus\bigoplus_{\alpha\in R^+}\operatorname{Lie}(G)^\alpha.
\]
\end{sgastatement}
\SGARefTarget{sga3:xxvi:thm:7:4}

\SGASubsection{7.4.1}{}
\SGARefTarget{sga3:xxvi:subsec:7:4:1}

Suppose first that the existence of the desired map
\(\alpha\mapsto\alpha^*\) has been proved.  By \SGARefLink{sga3:XXI:3.4.10}{Exposé XXI, 3.4.10},
\(s_\alpha\) is the unique element of \(W\) such that
\[
 s_\alpha(m)-m
\]
is a rational multiple of \(\alpha\) for every \(m\in M\).
Consequently \(s_\alpha\) is determined by \(\alpha\).  Since
\[
 \alpha^*(m)\alpha=m-s_\alpha(m),
\]
\footnote{\emph{Editor's note.} Here and again in \SGARefLink{sga3:xxvi:subsec:7:4:4}{7.4.4}, the
Polo--Gille edition restores the factor \(\alpha\) omitted after
\(\alpha^*(m)\) in the original text.}
the coroot \(\alpha^*\) is also determined by \(\alpha\), proving
the uniqueness of the map.

\SGASubsection{7.4.2}{}
\SGARefTarget{sga3:xxvi:subsec:7:4:2}

Let \(\alpha\in R\).  Denote by \(L_\alpha\), \(H_\alpha\), and
\(H_{-\alpha}\), respectively, the unique smooth subgroups of \(G\)
with connected fibers which contain \(L\) and satisfy
\textup{(cf.\ \SGARefLink{sga3:xxvi:proposition:6-1}{6.1}\textup{(i)})}
\[
\begin{aligned}
 \operatorname{Lie}(L_\alpha)
  &=\operatorname{Lie}(L)
    \oplus\bigoplus_{\gamma\in\mathbf Z\alpha\cap R}
       \operatorname{Lie}(G)^\gamma,\\
 \operatorname{Lie}(H_\alpha)
  &=\operatorname{Lie}(L)
    \oplus\bigoplus_{\gamma\in\mathbf N\alpha\cap R}
       \operatorname{Lie}(G)^\gamma,\\
 \operatorname{Lie}(H_{-\alpha})
  &=\operatorname{Lie}(L)
    \oplus\bigoplus_{\gamma\in-\mathbf N\alpha\cap R}
       \operatorname{Lie}(G)^\gamma.
\end{aligned}
\]
The group \(L_\alpha\) is a reductive subgroup of \(G\), and
\(H_\alpha,H_{-\alpha}\) are opposite parabolic subgroups of
\(L_\alpha\), with common Levi subgroup \(L\), by
\SGARefLink{sga3:xxvi:proposition:6-1}{6.1}\textup{(ii)}.

By \SGARefLink{sga3:xxvi:subsec:7:2}{7.2} there is therefore an element
\[
 s_\alpha\in\Norm_{L_\alpha}(Q)(S)/L(S)\subset W
\]
such that \(s_\alpha(H_\alpha)=H_{-\alpha}\).  One has
\[
 s_\alpha(\alpha)=-\alpha,
\]
because \(\alpha\), respectively \(-\alpha\), is the common divisor
of the elements of \(R\) occurring in \(H_\alpha\), respectively
\(H_{-\alpha}\).  Moreover \(s_\alpha^2=\mathrm{id}\), because
\(s_\alpha^2(H_\alpha)\) and \(H_\alpha\) are both opposite to
\(H_{-\alpha}\) relative to \(L\).  We have thus constructed an
element \(s_\alpha\in W\) with the properties
\[
\tag{x}\label{sga3:xxvi:subsection:7-4-2:equation:x}
 s_\alpha(\alpha)=-\alpha,\qquad s_\alpha^2=\mathrm{id},
\]
and
\[
\tag{xx}\label{sga3:xxvi:subsection:7-4-2:equation:xx}
 s_\alpha\text{ can be represented by an element of }L_\alpha(S).
\]
The construction is canonical in \(\alpha\); in particular,
\[
\tag{xxx}\label{sga3:xxvi:subsection:7-4-2:equation:xxx}
 w s_\alpha w^{-1}=s_{w(\alpha)}
 \qquad\text{for every }w\in W.
\]

\SGASubsection{7.4.3}{}
\SGARefTarget{sga3:xxvi:subsec:7:4:3}

We now prove
\[
\tag{xxxx}\label{sga3:xxvi:subsection:7-4-3:equation:xxxx}
 s_\alpha(m)-m\in\mathbf Z\alpha
 \qquad\text{for every }m\in M.
\]
Since \(S\) is connected, this assertion is local for the fpqc
topology.  We may therefore suppose that
\[
 L_\alpha=G_1
\]
is split relative to a maximal torus \(T_1\) of \(L\).  Let
\((G_1,T_1,M_1,R_1)\) be such a splitting.  The monomorphism
\(Q\to T_1\) identifies \(M\) with a quotient of \(M_1\); denote the
canonical map by
\[
 p:M_1\longrightarrow M.
\]

The image of \(R_1\) under \(p\), possibly including zero, consists
of the roots of \(G_1\) relative to \(Q\), hence of the elements of
\(R\) which are integral multiples of \(\alpha\).  Thus
\[
 p(R_1)\subset\mathbf Z\alpha.
\]
By \textup{\SGARefLink{sga3:xxvi:subsection:7-4-2:equation:xx}{(xx)}}, an element of \(\Norm_{G_1}(L)(S)\) induces
\(s_\alpha\) on \(Q\).  \SGARefLink{sga3:xxii:rem:5.10.10}{Exposé XXII, 5.10.10}, then gives a section
\[
 w\in(W_1)_S(S)
\]
which induces \(s_\alpha\) on \(Q\), where \(W_1\) is the Weyl group
of the root datum \((M_1,R_1,\ldots)\).  After restricting \(S\) if
necessary, we may suppose that \(w\in W_1\) and that
\[
 p\bigl(w(m_1)\bigr)=s_\alpha\bigl(p(m_1)\bigr)
 \qquad(m_1\in M_1).
\]
By definition, \(w\) is a product of reflections associated with
elements of \(R_1\).  Hence \(w(m_1)-m_1\) is an integral linear
combination of elements of \(R_1\).  It follows that
\[
 s_\alpha\bigl(p(m_1)\bigr)-p(m_1)
\]
is an integral linear combination of elements of
\(p(R_1)\subset\mathbf Z\alpha\), and is therefore an integral
multiple of \(\alpha\), as required.

\SGASubsection{7.4.4}{}
\SGARefTarget{sga3:xxvi:subsec:7:4:4}

We may now define \(\alpha^*\in M^*\) by
\[
 \alpha^*(m)\alpha=m-s_\alpha(m).
\]
By \textup{(x)},
\[
 \langle\alpha^*,\alpha\rangle=2.
\]
Moreover, \textup{\SGARefLink{sga3:xxvi:subsection:7-4-2:equation:xxx}{(xxx)}} implies that for every
\((\alpha,\alpha')\in R\times R\),
\[
 s_\alpha(\alpha')\in R,\qquad
 s_\alpha(\alpha'^*)=s_\alpha(\alpha')^*.
\]
By \SGARefLink{sga3:xxi:subsection:1-1}{Exposé XXI, 1.1}, the map
\(\alpha\mapsto\alpha^*\) therefore defines a root datum on
\((M,M^*)\).

\SGASubsection{7.4.5}{}
\SGARefTarget{sga3:xxvi:subsec:7:4:5}

Let \(W'\) be the Weyl group of this root datum, that is, the group
of transformations of \(M\) generated by the \(s_\alpha\).  Then
\(W'\subset W\).

Choose a total order on the free abelian group \(M\), and put
\[
 R^+=\{\alpha\in R\mid\alpha\geq0\}.
\]
This is a system of positive roots in \(R\).  Let \(w\in W\), and
represent it by an element
\[
 n\in\Norm_G(L)(S)=\Norm_G(Q)(S).
\]
Put \(P=H_{R^+}\), in the notation of \SGARefLink{sga3:xxvi:proposition:6-1}{6.1}.  Then \(P\) is a
parabolic subgroup of \(G\) with Levi subgroup \(L\), and
\[
 \Int(n)P=H_{w(R^+)}.
\]
It follows from \SGARefLink{sga3:xxvi:subsec:7:3}{7.3} that \(w(R^+)=R^+\) implies \(w=e\).  The group
\(W'\) acts transitively on the systems of positive roots in \(R\),
by \SGARefLink{sga3:XXI:3.3.7}{Exposé XXI, 3.3.7}, while the stabilizer of \(R^+\) in \(W\) is
trivial.  Therefore \(W=W'\).

It follows as well that \(W=W'\) acts simply transitively both on
the systems of positive roots in \(R\) and on the parabolic
subgroups of \(G\) having Levi subgroup \(L\).  This proves the last
assertion of \SGARefLink{sga3:xxvi:theorem:7-4}{7.4}.
\SGAQED
